% static: robustness of the loss-dominance conclusion (author: fill any TBD)
\begin{table}[t]
\centering
\caption{Robustness of the decomposition. The loss-dominance conclusion survives every modelling
choice examined; where a quantity is sensitive (last-mile loss magnitude, value of temperature
flexibility) it does not affect the ordering. Values are the loss main effect (\% of the gap) or
the copperplate regret sign, as applicable.}
\label{tab:robustness}
\begin{tabular}{lc}
\hline
Modelling choice & Loss dominance holds? \\
\hline
Remove the indefensible $\times$ multiplier (defensible $U$) & yes ($95.5\!\to\!95.9$\,\%) \\
Two lumped-loss calibrations (constant / heating-curve) & yes (regret $-0.5$/$-0.6$\,\%) \\
Demand-disaggregation rule (demand-share / uniform) & yes ($<0.3$\,pp) \\
Marginal-heat-cost assumption ($47$--$90$\,EUR/MWh) & yes (regret sign invariant) \\
Solver-gap tightening ($0.5\% \to 0.01\%$) & yes (loss unmoved; interaction resolves) \\
Flexible supply temperature (forward) & yes (saving is a loss effect) \\
\hline
\end{tabular}
\end{table}
